\usetheme{Madrid}
\usecolortheme{seahorse}
\setbeamertemplate{navigation symbols}{}
\setbeamertemplate{itemize items}[circle]
\setbeamertemplate{enumerate items}[default]

\usepackage{amsmath,amssymb,mathtools}
\usepackage{booktabs}

\newcommand{\OmegaSpace}{\Omega}
\newcommand{\F}{\mathcal{F}}
\newcommand{\B}{\mathcal{B}}
\newcommand{\R}{\mathbb{R}}
\newcommand{\Prob}{\mathbb{P}}
\newcommand{\E}{\mathbb{E}}
\newcommand{\Q}{\mathbb{Q}}
\newcommand{\one}{\mathbf{1}}
\newcommand{\vikiname}[1]{\par\smallskip{\scriptsize\color{blue!60!black}\textbf{LLMwiki:} \texttt{\detokenize{#1}}}}
\newcommand{\blankproof}{%
  \vspace{2mm}
  {\color{gray!55}\rule{\linewidth}{0.4pt}}\par\vspace{5mm}
  {\color{gray!55}\rule{\linewidth}{0.4pt}}\par\vspace{5mm}
  {\color{gray!55}\rule{\linewidth}{0.4pt}}\par\vspace{5mm}
  {\color{gray!55}\rule{\linewidth}{0.4pt}}\par
}
\newcommand{\proofblock}[1]{\ifteacher #1 \else \blankproof \fi}

\title[Chapter 1.2]{Chapter 1.2 Distributions}
\subtitle{Rick Durrett Probability: Theory and Examples}
\author{Hu Jingwu}
\date{}

\begin{document}

\begin{frame}
  \titlepage
  \vspace{-5mm}
  \begin{center}
    \small Built from the local Chapter 1 LLMwiki flash cards.
  \end{center}
\end{frame}

\begin{frame}{Roadmap for 1.2}
  \begin{block}{Learning goal}
    Move between laws, CDFs, densities/mass functions, atoms, and transformations.
  \end{block}
  \begin{enumerate}
    \item Distribution functions and probability laws.
    \item Atoms as jumps and countability of discontinuities.
    \item Normal tail estimates.
    \item Probability integral transform.
    \item Density transformations: monotone maps, lognormal, and squares.
  \end{enumerate}
\end{frame}

\begin{frame}{Distribution Function Properties}
  \begin{block}{Theorem}
    If $F(x)=\Prob(X\le x)$, then:
    \[
      F \text{ is nondecreasing},\qquad
      \lim_{x\to-\infty}F(x)=0,\qquad
      \lim_{x\to\infty}F(x)=1,
    \]
    and $F$ is right-continuous.
  \end{block}
  \begin{alertblock}{Pitfall}
    A CDF need not be left-continuous. Its jumps are atoms.
  \end{alertblock}
  \vikiname{Distribution function properties; ID: durrett_ch1_distribution_function}
\end{frame}

\begin{frame}{How to Prove CDF Properties}
  \begin{block}{Proof idea}
    Translate analytic properties of $F$ into event inclusions and continuity of probability.
  \end{block}
  \proofblock{
  \begin{proof}
    If $x\le y$, then $\{X\le x\}\subseteq\{X\le y\}$, so $F(x)\le F(y)$.
    For right-continuity, take $x_n\downarrow x$. Then
    \[
      \{X\le x_n\}\downarrow \{X\le x\}.
    \]
    By continuity from above,
    \[
      F(x_n)=\Prob(X\le x_n)\downarrow \Prob(X\le x)=F(x).
    \]
    The limits at $\pm\infty$ follow from $\{X\le n\}\uparrow\Omega$ and
    $\{X\le -n\}\downarrow\varnothing$.
  \end{proof}
  }
  \vikiname{Distribution function properties; ID: durrett_ch1_distribution_function}
\end{frame}

\begin{frame}{Every Valid CDF Determines a Law}
  \begin{block}{Theorem}
    Any function $F:\R\to[0,1]$ that is nondecreasing, right-continuous, and satisfies
    \[
      F(-\infty)=0,\qquad F(\infty)=1
    \]
    is the CDF of a unique probability measure on $\B(\R)$.
  \end{block}
  \begin{itemize}
    \item Define $\mu((a,b])=F(b)-F(a)$.
    \item Extend from half-open intervals to the Borel sigma-field.
  \end{itemize}
  \vikiname{Every valid CDF determines a law; ID: durrett_ch1_cdf_to_law}
\end{frame}

\begin{frame}{Atoms Are Jumps}
  \begin{block}{Fact}
    For a real random variable with CDF $F$,
    \[
      \Prob(X=x)=F(x)-F(x-),
      \qquad
      F(x-)=\lim_{y\uparrow x}F(y).
    \]
  \end{block}
  \proofblock{
  Since $\{X<x\}=\bigcup_{n\ge1}\{X\le x-1/n\}$, continuity from below gives
  \[
    \Prob(X<x)=\lim_{n\to\infty}F(x-1/n)=F(x-).
  \]
  Therefore
  \[
    \Prob(X=x)=\Prob(X\le x)-\Prob(X<x)=F(x)-F(x-).
  \]
  }
  \vikiname{Atoms are jumps of the CDF; ID: durrett_ch1_atoms_jumps}
\end{frame}

\begin{frame}{Standard Normal Tail Estimate}
  \begin{block}{Inequality}
    If $\chi$ is standard normal and $x>0$, then Durrett's normal-tail estimate gives
    \[
      (x^{-1}-x^{-3})e^{-x^2/2}
      \le
      \int_x^\infty e^{-y^2/2}\,dy
      \le
      x^{-1}e^{-x^2/2}.
    \]
  \end{block}
  Hence
  \[
    \Prob(\chi\ge x)=\frac{1}{\sqrt{2\pi}}\int_x^\infty e^{-y^2/2}\,dy.
  \]
  \vikiname{Standard normal tail estimate; ID: durrett_ch1_normal_tail_bound}
\end{frame}

\begin{frame}{Probability Integral Transform}
  \begin{block}{Fact}
    If $X$ has continuous CDF $F$, then
    \[
      F(X)\sim \operatorname{Unif}(0,1).
    \]
  \end{block}
  \begin{itemize}
    \item Proof route: compute $\Prob(F(X)\le y)$ through a quantile.
    \item Flat intervals require care, but continuity makes their probability contribution zero.
  \end{itemize}
  \vikiname{Probability integral transform; ID: durrett_ch1_probability_integral_transform}
\end{frame}

\begin{frame}{Monotone Density Change of Variables}
  \begin{block}{Theorem}
    Suppose $X$ has density $f$ on $[\alpha,\beta]$ and $g$ is strictly increasing and differentiable. If $Y=g(X)$, then
    \[
      f_Y(y)=
      \frac{f(g^{-1}(y))}{g'(g^{-1}(y))}
      \one_{\{g(\alpha)<y<g(\beta)\}}.
    \]
  \end{block}
  \proofblock{
  For $g(\alpha)<y<g(\beta)$,
  \[
    \Prob(Y\le y)=\Prob(X\le g^{-1}(y))
    =\int_{\alpha}^{g^{-1}(y)}f(x)\,dx.
  \]
  Differentiate with respect to $y$ and use
  $(g^{-1})'(y)=1/g'(g^{-1}(y))$.
  }
  \vikiname{Monotone density change of variables; ID: exercise_method_density_change_of_variables}
\end{frame}

\begin{frame}{Exercise 1.2.1}
  Let $X$ and $Y$ be random variables on $(\Omega,\F,P)$ and let
  $A\in\F$. Define $Z=X$ on $A$ and $Z=Y$ on $A^c$. Show that $Z$ is a
  random variable.
  \vikiname{Piecewise random variable measurability; ID: exercise_method_piecewise_random_variable_measurability}
\end{frame}

\begin{frame}{Exercise 1.2.1: Proof Skeleton}
  \proofblock{
  It is enough to check $\{Z\le x\}\in\F$ for each $x\in\R$. Write
  \[
    \{Z\le x\}
    =
    \bigl(A\cap\{X\le x\}\bigr)
    \cup
    \bigl(A^c\cap\{Y\le x\}\bigr).
  \]
  Each set on the right is measurable, so $Z$ is measurable.
  }
  \vikiname{Piecewise random variable measurability; ID: exercise_method_piecewise_random_variable_measurability}
\end{frame}

\begin{frame}{Exercise 1.2.2}
  Let $\chi$ be standard normal. Use Durrett's normal-tail estimate to
  find upper and lower bounds for $\Prob(\chi\ge4)$.
  \vikiname{Standard normal tail estimate; ID: durrett_ch1_normal_tail_bound}
\end{frame}

\begin{frame}{Exercise 1.2.2: Proof Skeleton}
  \proofblock{
  Durrett's estimate with $x=4$ gives
  \[
    \left(\frac14-\frac{1}{64}\right)e^{-8}
    \le
    \int_4^\infty e^{-y^2/2}\,dy
    \le
    \frac14e^{-8}.
  \]
  Therefore
  \[
    \frac{15}{64\sqrt{2\pi}}e^{-8}
    \le
    \Prob(\chi\ge4)
    \le
    \frac{1}{4\sqrt{2\pi}}e^{-8}.
  \]
  }
  \vikiname{Standard normal tail estimate; ID: durrett_ch1_normal_tail_bound}
\end{frame}

\begin{frame}{Exercise 1.2.3}
  Show that a distribution function has at most countably many
  discontinuities.
  \vikiname{Countability of CDF jumps; ID: exercise_method_countable_cdf_discontinuities}
\end{frame}

\begin{frame}{Exercise 1.2.3: Proof Skeleton}
  \proofblock{
  Let $\Delta F(x)=F(x)-F(x-)$. For $m,n\ge1$, set
  \[
    D_{m,n}=\{x\in[-n,n]:\Delta F(x)>1/m\}.
  \]
  Each $D_{m,n}$ is finite, since finitely many jumps in it have total size at most $1$ and each jump is bigger than $1/m$. Every discontinuity has $\Delta F(x)>0$, so it lies in some $D_{m,n}$. Hence the discontinuity set is a countable union of finite sets.
  }
  \vikiname{Countability of CDF jumps; ID: exercise_method_countable_cdf_discontinuities}
\end{frame}

\begin{frame}{Exercise 1.2.4}
  Suppose $F(x)=\Prob(X\le x)$ is continuous. Show that $Y=F(X)$ is
  uniform on $(0,1)$, i.e.
  \[
    \Prob(Y\le y)=y,\qquad 0\le y\le1.
  \]
  \vikiname{Probability integral transform; ID: durrett_ch1_probability_integral_transform}
\end{frame}

\begin{frame}{Exercise 1.2.4: Proof Skeleton}
  \proofblock{
  For $0\le y\le1$, use the generalized quantile
  \[
    x_y=\sup\{x:F(x)\le y\}.
  \]
  By monotonicity and continuity, $\{F(X)\le y\}$ has the same probability as a suitable half-line endpoint with CDF value $y$. Flat intervals carry no probability mass because $F$ is constant there. Thus
  \[
    \Prob(F(X)\le y)=y.
  \]
  }
  \vikiname{Probability integral transform; ID: durrett_ch1_probability_integral_transform}
\end{frame}

\begin{frame}{Exercise 1.2.5}
  Suppose $X$ has continuous density $f$ and is supported on
  $[\alpha,\beta]$. Let $g$ be strictly increasing and differentiable on
  $(\alpha,\beta)$. Find the density of $g(X)$.
  \vikiname{Monotone density change of variables; ID: exercise_method_density_change_of_variables}
\end{frame}

\begin{frame}{Exercise 1.2.5: Proof Skeleton}
  \proofblock{
  Let $Y=g(X)$. For $g(\alpha)<y<g(\beta)$,
  \[
    F_Y(y)=\Prob(g(X)\le y)=\Prob(X\le g^{-1}(y))
    =\int_{\alpha}^{g^{-1}(y)}f(x)\,dx.
  \]
  Differentiating,
  \[
    f_Y(y)=\frac{f(g^{-1}(y))}{g'(g^{-1}(y))},
  \]
  and $f_Y(y)=0$ outside $(g(\alpha),g(\beta))$.
  }
  \vikiname{Monotone density change of variables; ID: exercise_method_density_change_of_variables}
\end{frame}

\begin{frame}{Exercise 1.2.6}
  Suppose $X$ is normally distributed. Use the previous exercise to find
  the density of $\exp(X)$, the lognormal distribution.
  \vikiname{Monotone density change of variables; ID: exercise_method_density_change_of_variables}
\end{frame}

\begin{frame}{Exercise 1.2.6: Proof Skeleton}
  \proofblock{
  If $X\sim N(\mu,\sigma^2)$ and $Y=e^X$, then $g^{-1}(y)=\log y$ and
  $g'(g^{-1}(y))=y$. Hence, for $y>0$,
  \[
    f_Y(y)=\frac{1}{y\sigma\sqrt{2\pi}}
    \exp\left[-\frac{(\log y-\mu)^2}{2\sigma^2}\right],
  \]
  and $f_Y(y)=0$ for $y\le0$.
  }
  \vikiname{Monotone density change of variables; ID: exercise_method_density_change_of_variables}
\end{frame}

\begin{frame}{Exercise 1.2.7}
  \begin{enumerate}
    \item[(i)] If $X$ has density $f$, compute the distribution function and density of $X^2$.
    \item[(ii)] Specialize to $X$ standard normal to get the chi-square density with one degree of freedom.
  \end{enumerate}
  \vikiname{Density of a square transform; ID: exercise_method_square_transform_density}
\end{frame}

\begin{frame}{Exercise 1.2.7: Proof Skeleton}
  \proofblock{
  Let $Y=X^2$. For $y\ge0$,
  \[
    F_Y(y)=\Prob(-\sqrt y\le X\le\sqrt y)
    =\int_{-\sqrt y}^{\sqrt y}f(x)\,dx.
  \]
  Thus, for $y>0$,
  \[
    f_Y(y)=\frac{f(\sqrt y)+f(-\sqrt y)}{2\sqrt y}.
  \]
  If $X$ is standard normal, then
  \[
    f_Y(y)=\frac{1}{\sqrt{2\pi y}}e^{-y/2}\one_{(0,\infty)}(y).
  \]
  }
  \vikiname{Density of a square transform; ID: exercise_method_square_transform_density}
\end{frame}

\begin{frame}{Checklist: How to Attack 1.2 Problems}
  \begin{enumerate}
    \item Need to verify a CDF? Check monotonicity, right-continuity, and endpoint limits.
    \item Need an atom? Compute $F(x)-F(x-)$.
    \item Need countability of jumps? Group jumps by size and bounded interval.
    \item Need a transformed density? Compute the transformed CDF first, then differentiate.
    \item Need a normal tail? Keep the exponential term and constants separate.
  \end{enumerate}
  \vikiname{Chapter 1.2 flash cards and exercise method cards}
\end{frame}

\begin{frame}{Mini Practice}
  \begin{enumerate}
    \item If $F$ is continuous and strictly increasing, show $\Prob(F(X)\le y)=y$ directly using $F^{-1}$.
    \item If $X$ has density $f$, find the density of $aX+b$ for $a<0$.
    \item Show that a CDF can have only finitely many jumps larger than $0.01$.
    \item If $X\sim N(0,1)$, use the tail bound to estimate $\Prob(|X|\ge x)$.
  \end{enumerate}
\end{frame}

\end{document}
